% !TeX root = main.tex
% Synchronized with 00_摘要.md v0.6; Tsai format.
\begin{abstractzh}
\setcounter{page}{1}

移動機械臂於動態環境執行操作時，須同時處理障礙避讓、底盤與手臂運動耦合，以及控制命令與實際執行之間的差異。如何在運動限制下協調導航與操作，並兼顧追蹤品質及任務效率，為本研究關注的核心問題。

本研究提出以幾何預測與執行動態建模為核心的移動機械臂控制架構。導航部分以 \(SE(2)\) 建立位姿誤差模型，結合預測式控制屏障條件與掃描防護機制；協同操作部分將手臂位置設定點及追蹤動態納入預測，於有限時域內共同最佳化底盤與手臂輸入，協調末端追蹤與運動限制，並由視覺量測建立操作參考。

模擬結果顯示，所提出導航方法在 83 條共同評估路線中全部到達，其中 1 條發生接觸；MPPI 與 RPP 分別到達 64 條與 60 條。抽屜操作比較中，協同移動策略在符合既定任務判準下，相較於停車後操作策略縮短任務時間 13.38--14.05\%；五步預測相較單步預測，亦降低控制期間的末端位置誤差與手臂命令反轉率。上述結果顯示，所提出之預測與協同設計可於所測情境中兼顧操作品質與任務效率。

\keywordzh{\mbox{移動機械臂}、\allowbreak\mbox{幾何模型預測控制}、\allowbreak\mbox{全身協調控制}、\allowbreak\mbox{執行動態建模}、\allowbreak\mbox{控制屏障函數}、\allowbreak\mbox{移動中操作}}
\end{abstractzh}
